%%%%%%%%%%%%%%%%%
% This is an sample CV template created using altacv.cls
% (v1.7.1, 9 August 2023) written by LianTze Lim (liantze@gmail.com). Compiles with pdfLaTeX, XeLaTeX and LuaLaTeX.
%
%% It may be distributed and/or modified under the
%% conditions of the LaTeX Project Public License, either version 1.3
%% of this license or (at your option) any later version.
%% The latest version of this license is in
%%    http://www.latex-project.org/lppl.txt
%% and version 1.3 or later is part of all distributions of LaTeX
%% version 2003/12/01 or later.
%%%%%%%%%%%%%%%%

%% Use the "normalphoto" option if you want a normal photo instead of cropped to a circle
% \documentclass[10pt,a4paper,normalphoto]{altacv}

\documentclass[10pt,a4paper,ragged2e,withhyper]{altacv}
%% AltaCV uses the fontawesome5 and packages.
%% See http://texdoc.net/pkg/fontawesome5 and http://texdoc.net/pkg/academicons
%% for full list of symbols. You MUST compile with XeLaTeX or LuaLaTeX if you
%% want to use academicons.

% Change the page layout if you need to
\geometry{left=1.25cm,right=1.25cm,top=1.5cm,bottom=1.5cm,columnsep=1.2cm}

% The paracol package lets you typeset columns of text in parallel
\usepackage{paracol}

\usepackage{ctex}

% Change the font if you want to, depending on whether
% you're using pdflatex or xelatex/lualatex
% WHEN COMPILING WITH XELATEX PLEASE USE
% \setmainfont{Lato}

% Change the colours if you want to
\definecolor{SlateGrey}{HTML}{2E2E2E}
\definecolor{LightGrey}{HTML}{666666}
\definecolor{DarkPastelRed}{HTML}{450808}
\definecolor{PastelRed}{HTML}{8F0D0D}
\definecolor{GoldenEarth}{HTML}{E7D192}
\colorlet{name}{black}
\colorlet{tagline}{PastelRed}
\colorlet{heading}{DarkPastelRed}
\colorlet{headingrule}{GoldenEarth}
\colorlet{subheading}{PastelRed}
\colorlet{accent}{PastelRed}
\colorlet{emphasis}{SlateGrey}
\colorlet{body}{LightGrey}

% Change some fonts, if necessary
\renewcommand{\namefont}{\Huge\rmfamily\bfseries}
\renewcommand{\personalinfofont}{\footnotesize}
\renewcommand{\cvsectionfont}{\Large\rmfamily\bfseries}
\renewcommand{\cvsubsectionfont}{\large\bfseries}


% Change the bullets for itemize and rating marker
% for \cvskill if you want to
\renewcommand{\cvItemMarker}{{\small\textbullet}}
\renewcommand{\cvRatingMarker}{\faCircle}
% ...and the markers for the date/location for \cvevent
% \renewcommand{\cvDateMarker}{\faCalendar*[regular]}
% \renewcommand{\cvLocationMarker}{\faMapMarker*}


% If your CV/résumé is in a language other than English,
% then you probably want to change these so that when you
% copy-paste from the PDF or run pdftotext, you get
% sensible characters. But they can be changed here.
\ifpdf
  \input{glyphtounicode}
  \pdfglyphtounicode{f_f}{FB00}
  \pdfglyphtounicode{f_f_i}{FB03}
  \pdfglyphtounicode{f_f_l}{FB04}
  \pdfglyphtounicode{f_i}{FB01}
  \pdfgentounicode=1
\fi

%% Use (and optionally edit) this .tex if you
%% want to use an author-year reference style instead.
%% \input{pubs-authoryear.tex}

%% Use (and optionally edit) this .tex if you
%% want an originally numerical reference style.
%% \input{pubs-num.tex}

%% sample.bib contains your publications
%% \addbibresource{sample.bib}

\begin{document}
\name{马逸飞}
% \tagline{高级软件研发工程师 | 自动驾驶中间件与系统优化专家}
%% You can add multiple photos on the left or right
% \photoR{2.8cm}{Globe_High}
% \photoL{2.5cm}{Yacht_High,Suitcase_High}

\personalinfo{%
  % Not all of these are required!
  \email{yifeima98@gmail.com}
  \phone{(86) 15537560237}
  \location{上海, 中国}
  %\homepage{www.homepage.com}
  %\twitter{@twitterhandle}
  %\linkedin{your_id}
  \github{YIFEI-MA}
  %\orcid{0000-0000-0000-0000}
  %% You can add your own arbitrary detail with
  %% \printinfo{symbol}{detail}[optional hyperlink prefix]
  % \printinfo{\faPaw}{Hey ho!}[https://example.com/]

  %% Or you can declare your own field with
  %% \NewInfoFiled{fieldname}{symbol}[optional hyperlink prefix] and use it:
  % \NewInfoField{gitlab}{\faGitlab}[https://gitlab.com/]
  % \gitlab{your_id}
	%%
	%% For services and platforms like Mastodon where there isn't a
	%% straightforward relation between the user ID and the hyperlink,
	%% you can use \printinfo directly e.g.
	% \printinfo{\faMastodon}{@username@instace}[https://instance.url/@username]
	%% But if you absolutely want to create new dedicated info fields for
	%% such platforms, then use \NewInfoField* with a star:
	% \NewInfoField*{mastodon}{\faMastodon}
	%% then you can use \mastodon, with TWO arguments where the 2nd argument is
	%% the full hyperlink.
	% \mastodon{@username@instance}{https://instance.url/@username}
}

\makecvheader
%% Depending on your tastes, you may want to make fonts of itemize environments slightly smaller
% \AtBeginEnvironment{itemize}{\small}

%% Set the left/right column width ratio to 6:4.
\columnratio{0.63}

% Start a 2-column paracol. Both the left and right columns will automatically
% break across pages if things get too long.
\begin{paracol}{2}

\cvsection{工作经历}

\cvevent{软件研发工程师}{卡尔动力 (滴滴自动驾驶卡车业务)}{2021年9月 -- 至今}{上海}

\textbf{自研通信中间件开发}
\begin{itemize}
\item 深度参与设计高可用中间件架构，设计兼容DDS协议的通信层，支持混合组网与数据回放
\item 优化消息调度算法和内存管理，平均端到端延迟优化$\geq$20ms，P9999延迟降低$\geq$150ms
\item 主导10+核心节点迁移自研中间件，提升系统整体性能和可靠性
\item 开发优化消息队列、时间轮等核心组件，采用无锁队列和共享内存等技术，支撑节点的高性能高并发和高实时性需求
\end{itemize}

\textbf{分布式性能监控平台}
\begin{itemize}
\item 构建完善的在线和离线性能数据分析体系，实现30+核心指标可视化监控和分析
\item 集成Perfetto+eBPF+ftrace实现应用内到内核全链路追踪，性能瓶颈定位效率提升3倍
\item 建立多版本性能基线比对与自动告警机制，融合路测数据和仿真基线数据至CI/CD支撑高质量软件发布
\end{itemize}

\textbf{传感器驱动架构重构与优化}
\begin{itemize}
\item 实现兼容多型号LiDAR/Radar的统一驱动，并建立统一框架与实时健康监控告警系统和故障自恢复机制 
\item 通过无锁队列+零拷贝+SIMD优化，CPU占用率降低80\%，并通过切帧优化使数据延迟降低35ms
\end{itemize}

\textbf{高性能数据压缩与存储系统}
\begin{itemize}
\item 开发基于ZSTD流式压缩的PCAP采集系统，激光雷达数据体积压缩至ROS bag的17\%，大幅降低IO开销
\item 利用HEVC编码优化图像存储，All-intra下实现40\%体积减少
\item 构建压缩格式下数据回放仿真系统实现数据闭环，综合每年节省存储成本\$900万美金，支撑L4车队规模从50台扩展至300+台
\end{itemize}

\textbf{智能仿真与测试系统}
\begin{itemize}
\item 开发ViL/HiL混合仿真平台功能，实现20+自动化评价指标和本地仿真可视化工具链，支持故障注入等功能
\item 构建仿真一致性验证工具和确定性仿真框架，保障算法迭代质量
\end{itemize}

%\divider

%\cvevent{Product Engineer}{Google}{23 June 1999 -- 2001}{Palo Alto, CA}

%\begin{itemize}
%\item Joined the company as the first female hire
%\item Developed targeted advertisement in order to use user's search queries and show them related ads
%\end{itemize}

\cvsection{技术栈}

\cvtag{ROS/DDS}
\cvtag{gRPC/Protobuf}
% \cvtag{实时系统优化}
\cvtag{C++}
\cvtag{Python}
\cvtag{CUDA/并行计算}
\cvtag{eBPF/ftrace}
\cvtag{Docker}
\cvtag{分布式系统}
\cvtag{QT}
\cvtag{PyTorch}
\cvtag{OpenCV}

\switchcolumn

\cvsection{教育背景}

\cvevent{计算机科学荣誉学士学位}{阿尔伯塔大学}{2017.09 -- 2021.06}{加拿大}
\begin{itemize}
\item 成绩：前10\% （ Dean's Honor Roll）
\item 核心课程：操作系统、计算机视觉、机器学习、自然语言处理、分布式系统
\item 项目经历：视频运动物体分割，基于MCTS的围棋AI，安卓打车应用
\end{itemize}

\cvsection{语言能力}

\cvskill{中文}{5}
\cvskill{英语}{4}

\cvsection{兴趣爱好}

\cvtag{风光摄影}
\cvtag{围棋（三段）}
\cvtag{篮球}

\cvsection{项目经历}

\cvevent{低延迟视频推流系统}{E-Stop核心功能}{2024.06 -- 2024.09}{}
\begin{itemize}
\item 在Orin平台开发多路4K视频处理流水线（去畸变→拼接→RTSP推流）
\item 通过CUDA零拷贝+NVENC（AV1/HEVC）硬编码实现低开销（CPU单核<40\%）
\item 低带宽Radio通信场景下的实时视频传输GTG延迟$\leq$200ms
\end{itemize}

\cvevent{自动驾驶仿真Smart Agent}{清华大学车辆工程学院}{2022.09 -- 2023.06}{}
\begin{itemize}
\item 合作开发Virtual Sim中的智能交通参与者系统
\item 实现基于规则、社会力模型和强化学习的混合多向避让和路径规划算法
\item 提升仿真场景真实性和可交互性，支撑大规模算法验证与测试
\end{itemize}

%\divider

%\cvevent{Bug Bounty Hunter}{Acme Corp}{Spring 2017 -- present}{Remote}
%\begin{itemize}
%\item Achieved over \$5,000 in bug bounties
%\item Reported several critical vulnerabilities
%\end{itemize}

%% Switch to the right column. This will now automatically move to the second
%% page if the content is too long.

% \cvsection{技术能力矩阵}

% \begin{tabular}{@{}ll@{}}
%   \textbf{核心领域} & 中间件架构 \textbullet\ 性能优化 \textbullet\ 实时系统 \textbullet\ 仿真平台 \\
%   \textbf{编程语言} & C++20/STL \textbullet\ Python3 \textbullet\ CUDA \textbullet\ Bash \\
%   \textbf{系统能力} & Linux内核调优 \textbullet\ eBPF/Perfetto \textbullet\ 无锁编程 \textbullet\ DDS/ROS2 \\
%   \textbf{工具链} & NVIDIA Jetson \textbullet\ ZSTD/HEVC \textbullet\ Docker/K8s \textbullet\ CI/CD \\
%   \textbf{算法基础} & 多传感器融合 \textbullet\ 强化学习 \textbullet\ 分布式系统 \textbullet\ 数据压缩 \\
% \end{tabular}

% \cvsection{核心技能}

% \cvskill{C++/系统编程}{5}
% \cvskill{Python/脚本开发}{5}
% \cvskill{CUDA/GPU编程}{4}
% \cvskill{Linux内核/eBPF}{4}
% \cvskill{ROS/DDS中间件}{5}
% \cvskill{性能优化/分析}{5}
% \cvskill{分布式系统}{4}
% \cvskill{机器学习基础}{3}



% \cvsection{教育背景}

% \cvevent{计算机科学荣誉学士学位}{阿尔伯塔大学}{2017.09 -- 2021.06}{加拿大}
% \begin{itemize}
% \item 成绩排名：前10\%
% \item 核心课程：操作系统、计算机视觉、机器学习、自然语言处理
% \item 项目经历：视频运动物体分割，基于蒙特卡洛树搜索的围棋AI，安卓打车应用
% \end{itemize}

%\divider

%\cvevent{M.Sc.\ in Your Discipline}{Your University}{Sept 2001 -- June 2003}{}

% \cvsection{核心优势}

% \begin{quote}
% ``具备深厚的自动驾驶系统架构经验，专精于高性能中间件开发与系统优化。拥有完整的产品开发生命周期经验，从底层驱动到上层应用全栈技术能力。''
% \end{quote}

% \cvsection{技术亮点}

% \cvachievement{\faTrophy}{成本优化}{年度节省存储成本\$900万美金}

% \cvachievement{\faRocket}{性能提升}{CPU占用率降低80\%，延迟优化150ms}

% \cvachievement{\faCogs}{系统架构}{设计高并发无锁数据处理架构}

% \cvachievement{\faChartLine}{监控平台}{构建30+指标的分布式监控系统}

%% Yeah I didn't spend too much time making all the
%% spacing consistent... sorry. Use \smallskip, \medskip,
%% \bigskip, \vspace etc to make adjustments.
\medskip

%% Yeah I didn't spend too much time making all the
%% spacing consistent... sorry. Use \smallskip, \medskip,
%% \bigskip, \vspace etc to make adjustments.
% \medskip

%% \cvsection{Referees}

%% \cvref{name}{email}{mailing address}
%% \cvref{Prof.\ Alpha Beta}{Institute}{a.beta@university.edu}
%% {Address Line 1\\Address line 2}

%% \cvref{Prof.\ Gamma Delta}{Institute}{g.delta@university.edu}
%% {Address Line 1\\Address line 2}

\end{paracol}


\end{document}